% ============================================================
% Introduction
% ============================================================

A growing body of theory predicts that \textit{anticipated} third-party
intervention shapes civil war onset
\citep{cetinyan2002,cunningham2016,dudley2024,gibilisco2022,kuperman2008,kyddstraus2013,lango2023,sambanis2020},
yet the quantity has never been measured systematically.  No existing measure
covers the full population of potential interveners, varies annually with
shifting alliances and foreign policy alignments, and treats intervention
direction---government-biased versus opposition-biased---as the primary
quantity of interest.  This paper constructs such a measure and uses it
to test the hypothesis that the shadow of intervention deters and
emboldens.

The construction is the core contribution.  A machine-learning ensemble
trained on the \citet{regan2000} dataset of military
interventions (1946--2014) predicts, for each directed dyad in a given
year, the probability and direction of intervention.  Dyad-level
predictions are aggregated across all potential interveners to produce a
country-year shadow, disciplined by a Nash fixed-point condition
requiring each state's predicted probability to be self-consistent as an
input to the others' predictions.  The resulting measures are tested in
a standard country-year onset framework \citep{fearon2003}: expected
government-biased intervention enters negatively (deterrence), expected
opposition-biased intervention enters positively (emboldening), with the
directional pattern robust across specifications, measurement draws, and
a country fixed-effects estimator.

The closest existing approaches each sacrifice scope for
identification.  \citet{cunningham2016} uses the Lake
security hierarchy as a proxy for anticipated government-biased
intervention---a static, unidirectional measure restricted to the US
patron-client channel.  The present paper generalizes in three
directions: it covers all potential interveners, both sides of
intervention, and derives probabilities from a calibrated classifier
rather than structural position.
\citet{gibilisco2022} derive equilibrium intervention
expectations from a structural game among the P5, gaining clean
identification at the cost of restricting the player set to a size that
keeps joint-game estimation tractable.  In the Regan data, P5 states
account for fewer than half of all military interventions; the remaining
53 intervening states---neighbors, coethnics, regional rivals, Cold War
proxies---represent the majority of actual deployments and are driven by
logics qualitatively different from the global-order interests that
animate P5 behavior.  \citet{lango2023} develops a
formal model in which the threat of rebel-sided intervention can
simultaneously deter onset---by compelling governments to concede---and
encourage it---by raising rebels' expected return from fighting; the
utility parameters in the present framework can represent these
nonmonotone effects.

This paper takes a learned-proxy approach: a flexible first-stage model
constructs the latent quantity, and a parametric second-stage model tests
the theory.  This two-stage design is an increasingly common way to
bring machine learning into social science inference
\citep{knoxlucascho2022}, and arguably the least controversial
\citep{moruccispirling2024}: the ensemble does the measurement work it is
suited for, while the onset regression retains the interpretability that
causal claims require.  The design nests the structural functional
form---the unpenalized multinomial logit in the ensemble library directly
corresponds to the softmax best-response function of
\citet{gibilisco2022}---so whether that parametric
assumption fits the data is answered empirically rather than imposed by
construction.  (It earns less than one percent of the ensemble weight;
the data strongly prefer nonparametric alternatives.)

The learned-proxy approach has well-documented pitfalls.
\citet{knoxlucascho2022} catalog these systematically:
measurement-stage uncertainty is routinely ignored, proxy quality is
asserted rather than tested, and the disconnect between prediction and
inference is papered over.  This paper follows their recommendations
closely: Stage~1 performance is validated out-of-fold across 25
imputation draws; measurement uncertainty is propagated into Stage~2 via
a two-stage pairs cluster bootstrap that averages across all 25 draws;
calibration is verified by reliability diagrams; and a labeled-only
robustness check confirms that the onset estimates are not driven by
out-of-sample extrapolation.  The honest accounting of generated-regressor
uncertainty is itself a result: the corrected standard errors are roughly
2.9--3.5 times larger than naive MLE output, with a variance decomposition
showing that approximately 88--92\% of the total coefficient variance originates
in the measurement stage and only 8--12\% from primary-analysis sampling.  The
measurement model, not the regression, is the binding source of
uncertainty---precisely the situation \citet{knoxlucascho2022} warn about.
The two shadow variables are correlated at $r = 0.94$, which inflates
individual confidence intervals, but tested jointly they reject the
baseline at $p < 0.001$; the directional pattern holds across all 25
measurement draws, ten specifications, and country fixed effects.
The shadow subsumes both the Lake security hierarchy
\citep{cunningham2016} and the \citet{gibilisco2022} structural P5
estimates: neither adds explanatory power once the shadow variables are
included.

The paper proceeds as follows.  Section~\ref{sec:motivation} develops the theoretical
framework linking intervention expectations to the onset decision.
Section~\ref{sec:constructing} describes the construction and validation of the
shadow measure.  Section~\ref{sec:decision} presents the onset analysis, including
measurement-corrected inference and heterogeneous-utility extensions.
